%%%PAGE 119 rot=0 legibility=good summary=物块落在弹簧上(A/O/B分析、v-t图、三种释放)、三类临界条件、三球杆连接体两题(机械能/动量)与三球V形杆模型
%%%BLOCK id=119-01 ch=03 sec=弹簧·物块落在弹簧上 kind=knowledge alt=06 cont=none y=0.07-0.21
% 标题“物块”二字前有一个难辨的小符号，未转录
\textbf{\underline{物块落在弹簧上}}

%FIG p119_a 0.075 0.10 0.145 0.188
\notefig[0.25]{figures/p119_a.png}
% 图中标注：弹簧上方一个圆圈(物块)，弹簧旁自上而下标 A、O、B 三个位置

\begin{itemize}
\item[A] $kx<mg$\quad $v\uparrow$\quad $a$向下$\downarrow$
\item[O] $kx=mg$\quad $v_{\max}$
\item[B] $kx>mg$\quad \struck{\unclear{…}}\ $v\uparrow$\quad $a$向上$\downarrow$ % CHECK: B 行 v、a 箭头方向与物理预期相反(预期 v 减小、a 向上增大)；kx>mg 后另有一处被涂改的字
\end{itemize}

$v_{\max}$不在刚碰处

%FIG p119_b 0.405 0.121 0.592 0.198
\notefig[0.25]{figures/p119_b.png}
% 图中标注：纵轴 v，$V_m$ 等标注，曲线上一点 D，横轴 t

\[
\left\{\begin{aligned}
&\text{有}h\text{释放}\quad a_{\text{下}}>g\\
&\text{无}h\text{释放}\quad a_{\text{下}}=g\\
&\text{提着释放}\quad a_{\text{下}}<g
\end{aligned}\right.
\qquad\text{动能定理证明}
\]
%%%END
%%%BLOCK id=119-02 ch=03 sec=动力学临界问题 kind=knowledge alt=none cont=none y=0.21-0.28
变速临界\quad $a=0$\ $v_{\max}$\quad $v=0$\ $a_{\max}$

分离临界\quad $N=0$\ $a_1=a_2$\ $v_1=v_2$\quad 不分离\ $N\ge 0$

滑动临界、整体隔离法
%%%END
%%%BLOCK id=119-03 ch=07 sec=杆绳连接体·机械能与动量 kind=problem alt=04,07 cont=none y=0.29-0.62
\begin{problem}[1.]
%FIG p119_c 0.065 0.30 0.255 0.405
\notefig[0.25]{figures/p119_c.png}
% 图中标注：m、L、m（上方）；A、B；L、L；m、C

① 求AB相碰瞬间速度

② AB速度与C到细杆距离的关系
\begin{solution}
动能定理

① \[ mg\left(L-\frac{\sqrt{3}}{2}L\right)=\frac{1}{2}\cdot 2mv^2 \qquad v=\sqrt{\left(1-\frac{\sqrt{3}}{2}\right)gL} \]

水平方向动量守恒\quad $mv_A+mv_B=0$

%FIG p119_d 0.052 0.445 0.15 0.535
\notefig[0.25]{figures/p119_d.png}
% 图中标注：$\theta$、$L$、$h$

② \[ mg\left(h-\frac{\sqrt{3}}{2}L\right)=\frac{1}{2}\cdot 2mv^2+\frac{1}{2}mv_C^2 \]
\[ v\cos\theta=v_C\sin\theta \]
\[ \tan\theta=\frac{h}{\sqrt{L^2-h^2}}\qquad\therefore\ v=\sqrt{\frac{gh^2\left(2h-\sqrt{3}L\right)}{L^2+h^2}} \]
\end{solution}
\end{problem}
%%%END
%%%BLOCK id=119-04 ch=07 sec=动量与能量综合·多球系统 kind=problem alt=06 cont=none y=0.63-0.79
\begin{problem}[2.]
%FIG p119_e 0.085 0.642 0.265 0.722
\notefig[0.25]{figures/p119_e.png}
% 图中标注：$m_1$、$M$、$m_2$（水平连线），$M$ 处向上箭头 $v_0$

求\ 1\ 2\ 相遇时速度
\begin{solution}
\begin{align*}
Mv_0&=(M+2m)v_y\\
\frac{1}{2}Mv_0^2&=\frac{1}{2}Mv_y^2+\frac{1}{2}\,2m\left(v_x^2+v_y^2\right)\\
v_x&=\sqrt{\frac{M}{M+2m}}\,v_0\qquad v_y=\frac{Mv_0}{M+2m}\\
v&=\sqrt{v_x^2+v_y^2}=\sqrt{\frac{2M(M+m)}{M+2m}}\,v_0 % CHECK: 根号是否只盖分子；若按 v_x、v_y 代入，分母应为 (M+2m) 而非根号下的 M+2m
\end{align*}
\end{solution}
\end{problem}
%%%END
%%%BLOCK id=119-05 ch=07 sec=杆连三球·V形模型 kind=problem alt=07,03 cont=none y=0.80-1.00
同
%FIG p119_g 0.535 0.795 0.945 0.86
\notefig[0.4]{figures/p119_g.png}
% 图中为四幅小示意图（竖直两球杆、倒V形、水平三球带向上箭头、斜靠直角装置），“同”字在其左侧

%FIG p119_f 0.005 0.805 0.215 0.962
\notefig[0.25]{figures/p119_f.png}
% 图中标注：A、m（顶点）；L、L；B、C 各 m

bc：$v$\ 0\ 有\ 0\quad $W_{\text{正}}$\ $W_{\text{负}}$

$F_{\text{杆}}=0$时\ $a_A=g$\ 其余时刻小于$g$

$v_{A\text{地}}=\sqrt{2gL}$\quad $mgL=\frac{1}{2}mv^2$\quad $v_B=v_C=0$\ 水平动量抵消

A机械能最小时\quad $v_A=0$\quad $F_{\text{杆}}=0$\quad $v_b=v_c=\dfrac{\sqrt{2gL}}{2}$\quad bc对地压力为$mg$

\unclear{有超\ 失重} % 页面底边被截断，字迹不全
%%%END
%%%PAGEEND 119
